\documentclass[11pt]{article}
\usepackage{amsmath,amssymb,amsthm,fullpage}
\newtheorem{proposition}{Proposition}
\title{An Audit of Recursive Exterior Cauchy Interpolation}
\author{}
\date{November 23, 2026}
\begin{document}
\maketitle

\paragraph{Status.}
This note tests a specific attempted implementation of the butterfly
rectangle decoder with label-dependent signed residues. It gives a
complete finite exact audit and a written obstruction for four
coalescing poles. Neither is a lower bound for unrestricted
HeadComplexity. They do not prove either accepted asymptotic target.
The original model is unchanged, and these results are not Lean checked.

\section{The recursive interpolation proposal}
At each right child in an output-prefix binary tree, take as query
profile the indicator of that child's leaves. If the child has $M$
leaves, give it $M$ exterior Cauchy queries. A permutation assigns
the poles to leaves, and all signed residues may depend on the entire
desired sign table. A subtree in isolation is a square Cauchy
interpolation problem, but the same residues must satisfy all its
descendant constraints.

The first useful augmented test has four poles, four root queries,
two queries for one two-leaf child, and two singleton queries. It
therefore asks for eight independent signs from four residues and
one of $24$ pole permutations. The exact results below show that the
most direct compact symmetric choice misses some tables.

\section{Complete finite exact test}
Put
\[
 \xi_j=10^{-5}+10^{-7}a_j,\qquad
 (a_0,a_1,a_2,a_3)=(-3,-1,1,3).
\]
For each $\alpha\in\{1/2,2,4\}$, the eight query times are
\[
 t_i=99/100+10^{-3}u_i,
 \qquad
 (u_0,\ldots,u_7)=(-3,-1,1,3,-\alpha,\alpha,0,0).
\]
The output leaf sets are
\[
 S_0=S_1=S_2=S_3=\{0,1,2,3\},\quad
 S_4=S_5=\{1,3\},\quad S_6=\{2\},\quad S_7=\{3\}.
\]
For a permutation $\pi$ and arbitrary real residues $c_j$, evaluate
\[
 f_i=\sum_{j:\pi(j)\in S_i}\frac{c_j}{(t_i-\xi_j)^2}.
 \tag{1}\label{eq:test}
\]
The companion checker proves that, for each of the three choices of
$\alpha$, the union over all $24$ permutations realizes exactly
$240$ of the $256$ strict sign tables. Each single permutation
realizes $116$ tables. Encode a positive sign in row $i$ by bit $i$
of an integer. The missing integers are exactly
\[
 21,26,37,42,85,90,101,106,
 149,154,165,170,213,218,229,234.
\]
Equivalently, the root's four signs alternate and the two-leaf
child's two signs alternate; the singleton signs are arbitrary.

This is complete exact verification despite degeneracies in the
central hyperplane arrangements. For every permutation and every
claimed realized table, the checker evaluates a rational residue
witness in the original matrix and checks all eight strict signs.
For every other table it checks a nonzero row dependence whose
coefficients have the requested signs on its support. The weighted
sum of hypothetical scores with those strict signs would be nonzero,
contradicting that dependence. Thus both presence and absence are
certified by rational arithmetic. No floating-point infeasibility
report is used.

\section{A uniform obstruction for four coalescing poles}
The following written statement explains the finite test and allows
the residues to vary arbitrarily with the pole separation.

\begin{proposition}\label{prop:uniform}
Let $p\geq1$ be an integer and $J=[a,b]\subset(0,\infty)$ satisfy
\[
 \frac{pb}{p+2}<a.
 \tag{2}\label{eq:interval}
\]
Fix four distinct ordered points $q_0<q_1<q_2<q_3$ in $J$ and
the four numbers $a_j=(-3,-1,1,3)_j$. There exists $\epsilon_0>0$
such that, whenever $0<\epsilon<\epsilon_0$ and
\[
 F_\epsilon(t)=\sum_{j=0}^3\frac{c_j}{(t-\epsilon a_j)^p}
\]
has strictly alternating signs at the four $q_i$, every two-term
partial sum of $F_\epsilon$ is nonzero and has constant sign on $J$.
\end{proposition}
\begin{proof}
Suppose the assertion fails along a sequence $\epsilon\downarrow0$.
Pass to a subsequence on which the same pair of terms fails, and
normalize the residues by a common positive factor so that their
first four scaled moments
\[
 m_k=\sum_j c_j(\epsilon a_j)^k,\quad 0\leq k\leq3,
 \qquad \max_k|m_k|=1
\]
have this normalization. The moments are not all zero because the
four-node Vandermonde matrix is invertible. Passing to a further
subsequence makes the moment vector converge to a vector of maximum
norm one. Denote its limiting coordinates again by $m_k$.

Vandermonde inversion gives $c_j=O(\epsilon^{-3})$. Taylor expansion
on the compact positive interval $J$ therefore yields uniform
convergence
\[
 F_\epsilon(t)\longrightarrow
 \sum_{k=0}^3\binom{p+k-1}{k}\frac{m_k}{t^{p+k}}
 =\frac{P(t)}{t^{p+3}},
\]
where $P$ is a nonzero polynomial of degree at most three.
The strict alternation supplies three real zeros between consecutive
$q_i$. It follows that $P$ has three roots $r_1,r_2,r_3$ in $J$,
counted with multiplicity. One elementary justification also handles
zeros tending to a common endpoint: multiply $F_\epsilon$ by its
positive common denominator $\prod_j(t-\epsilon a_j)^p$.
The resulting polynomials converge coefficientwise to
$t^{3p-3}P(t)$. Divide each by the three linear factors furnished
by the three zeros, pass to convergent subsequences of those zeros,
and use continuity of monic polynomial division. Their limiting
product divides $t^{3p-3}P(t)$, hence divides $P$ since $J$ excludes
zero.

Consequently
\[
 P(t)=A(t-r_1)(t-r_2)(t-r_3),\qquad A\ne0.
\]
In particular $m_2,m_3$ are nonzero, and comparison of coefficients
gives
\[
 -\frac{p m_3}{m_2}
 =\frac{3p}{p+2}\frac{r_1r_2r_3}
 {r_1r_2+r_1r_3+r_2r_3}.
 \tag{3}\label{eq:root}
\]
The right side is $p/(p+2)$ times the harmonic mean of the three
roots, and is therefore at most $pb/(p+2)<a$.

For clarity, the two relevant columns of the inverse Vandermonde
matrix are
\[
 w=\frac1{48}(-1,3,-3,1),\qquad
 u=\frac1{16}(1,-1,-1,1).
\]
They have moments $\sum_j w_j a_j^k=\mathbf1_{k=3}$ and
$\sum_j u_j a_j^k=\mathbf1_{k=2}$ for $0\leq k\leq3$.
Thus the leading residue terms are
$\epsilon^{-3}m_3w+\epsilon^{-2}m_2u$.
For any nonsymmetric pair of nodes, the sum of its two $w$
coordinates is nonzero. Multiplying that pair's partial score by
$\epsilon^3$ gives uniform convergence to a nonzero constant
multiple of $m_3/t^p$. Its sign is therefore constant on $J$ for
all sufficiently small $\epsilon$.

For the outer symmetric pair, the two $w$ coordinates sum to zero,
while both $\sum u_j$ and $\sum a_jw_j$ on that pair equal $1/8$.
For the inner symmetric pair both sums equal $-1/8$. Multiplication
of either partial score by $\epsilon^2$ therefore gives the uniform
limit
\[
 \pm\frac18\left(\frac{m_2}{t^p}
                  +\frac{p m_3}{t^{p+1}}\right).
\]
Its sole positive candidate zero is \eqref{eq:root}, which lies
strictly below $J$ by \eqref{eq:interval}. This limit is again
bounded away from zero on $J$. Every pair consequently has constant
nonzero sign for all sufficiently late terms of the supposed
counterexample sequence, a contradiction.
\end{proof}

For $p=2$ and three coalescing root zeros at $T$, the symmetric-pair
zero is $T/2$. For $p=1$ it is $T/3$. The same proof applies after
adding a common positive offset to all poles and query times.
This identifies the failure of the plain subtree-indicator proposal.
It does not address other color profiles, label-dependent query
designs, anchor cancellation rules, or arbitrary original-model heads.

\end{document}
